% ID: SEL-D0050-OB-E04
% ESTADO: validada. S051.
\textbf{a)} El intervalo normal con $\sigma$ conocida tiene centro $\bar x$ y semiamplitud $E=z_{1-\alpha/2}\sigma/\sqrt n$. Los extremos impresos dan
\[
\bar x=\tfrac{188{,}18+208{,}82}{2}=\boxed{198{,}50\text{ min}},
\qquad E=\tfrac{208{,}82-188{,}18}{2}=10{,}32\text{ min}.
\]
Con $z_{0{,}995}\approx2{,}57583$, al despejar se obtiene $n=(75z_{0{,}995}/10{,}32)^2\approx350{,}43$. La reconstrucción aproximada da $\boxed{n\approx350}$. De hecho, si se usa el cuantil aproximado $2{,}575$ y $n=350$, los extremos redondean a los impresos.

Con el cuantil exacto, $n=350$ da $E\approx10{,}3263$ y $n=351$ da $E\approx10{,}3116$. Ningún tamaño entero da un margen entre $10{,}315$ y $10{,}325$, como requerirían ambos extremos redondeados a centésimas con centro $198{,}50$. Por eso el tamaño se presenta como reconstrucción aproximada.

\textbf{b)} Al $96\%$, $z_{0{,}98}\approx2{,}05375$. El margen pedido para $n=500$ es
\[
E=2{,}05375\,\frac{75}{\sqrt{500}}\approx\boxed{6{,}89\text{ min}}.
\]
